\chapter{String Diagrams and Graphical Affine Algebra}\label{sec:gaa}

Having established the categorical background, we now introduce string diagrams
and the calculi of Graphical Linear Algebra and Graphical Affine Algebra
\cite{bonchi2019affine}. Readers familiar with these calculi may use this chapter
as a reference for the notation and conventions adopted here.

\section{String diagrams}\label{sec:string-diagrams}

A symmetric monoidal category has a sound two-dimensional notation: its
\emph{string diagrams}~\cite{joyal1991geometry,selinger2010survey}.  A wire
represents an object and a box with input and output wires represents a
morphism.  Reading from left to right, joining the outputs of $f$ to the inputs
of $g$ represents $g\circ f$; placing $f$ above $g$ represents $f\otimes g$.
A bare wire is an identity and crossing wires represents symmetry.  The
interchange law is built into the geometry: disjoint boxes may move past one
another because their relative horizontal order carries no information.
Consequently, diagrams are considered equal up to deformations that preserve
connectivity, boundary, and the order of ports on each box.  The basic
translation between categorical notation and string diagrams is displayed in
Figure~\ref{fig:appendix-string-diagram-dictionary}.

\begin{figure}[H]
\centering
\begin{tikzpicture}[axpic]
  % Identity and a general morphism
  \draw (-0.6,1.1) -- (0.6,1.1);
  \node at (0,0.45) {$\mathrm{id}_A:A\to A$};

  \pic[val=f] (single) at (2.6,1.1) {diagram};
  \node at (2.6,0.45) {$f:A\to B$};

  % Sequential composition
  \pic[val=f] (cf) at (4.6,1.1) {diagram};
  \pic[val=g] (cg) at (5.9,1.1) {diagram};
  \draw (cf-out-0) -- (cg-in-0);
  \node at (5.25,0.45) {$g\circ f$};

  % Tensor product
  \pic[val=f] (tf) at (7.8,1.65) {diagram};
  \pic[val=g] (tg) at (7.8,0.75) {diagram};
  \node at (7.8,0.1) {$f\otimes g$};

  % Symmetry
    \draw (9.5,1.45) to[out=0,in=180] (10.7,0.75);
    \draw (9.5,0.75) to[out=0,in=180] (10.7,1.45);
  \node at (10.1,0.1) {$\sigma_{A,B}$};
\end{tikzpicture}
\caption{String-diagram notation for the structure of a symmetric monoidal
category.  Diagrams are read from left to right: horizontal connection is
composition, while vertical juxtaposition is the tensor product.}
\label{fig:appendix-string-diagram-dictionary}
\end{figure}

The labels beneath the pictures record the corresponding morphisms.  In
particular, the two wires in the drawing of $f\otimes g$ do not interact: they
represent processes performed in parallel.  The crossing represents the
canonical reordering $A\otimes B\to B\otimes A$, not an additional operation
chosen for the particular objects.

\begin{figure}[H]
    \centering
\begin{tikzpicture}[axpic]
  \pic[val=f] (f) at (0,0.7) {diagram};
  \pic[val=g] (g) at (1.5,0.7) {diagram};
  \draw (f-out-0) -- (g-in-0);
  \pic[val=h] (h) at (0,-0.7) {diagram};
  \pic[val=k] (k) at (1.5,-0.7) {diagram};
  \draw (h-out-0) -- (k-in-0);
  \node at (3.0,0) {$=$};
  \pic[val=f] (f2) at (4.0,0.7) {diagram};
  \pic[val=g] (g2) at (5.5,0.7) {diagram};
  \draw (f2-out-0) -- (g2-in-0);
  \pic[val=h] (h2) at (4.0,-0.7) {diagram};
  \pic[val=k] (k2) at (5.5,-0.7) {diagram};
  \draw (h2-out-0) -- (k2-in-0);
  \node[below=2 of f-node] {$(g\circ f)\otimes(k\circ h)$};
  \node[below=2 of f2-node] {$(g\otimes k)\circ(f\otimes h)$};
\end{tikzpicture}
\caption{The interchange law expressed as a string diagram.}
\label{fig:appendix-interchange}
\end{figure}

For $\mathbf{Rel}$ the monoidal product is Cartesian product on objects.  If
$R:X\to Y$ and $S:U\to V$, then
\[
 R\otimes S:X\times U\longrightarrow Y\times V,
 \qquad
 ((x,u),(y,v))\in R\otimes S
 \Longleftrightarrow xRy\text{ and }uSv.
\]
The singleton set is the unit.  Sequential composition eliminates shared
variables; monoidal composition juxtaposes independent systems.

String diagrams encode structural bookkeeping, such as the interchange law
and the associativity of composition, in their geometry. This allows derivations
to focus on the operations and constraints of the system rather than repeatedly
invoking structural identities.

Diagrammatic equations can also be applied locally: a subdiagram may be replaced
by another with the same boundary without changing the surrounding diagram.
For the optimization language developed later, sequential composition denotes
the infimum over shared variables. The notation therefore allows components
to be manipulated without repeatedly expanding the corresponding infima.

\section{Frobenius structure and hypergraph categories}\label{sec:main-hypergraph}

The resulting categories are symmetric monoidal, and it is precisely this
monoidal backbone that will provide the algebraic structure in which we
are interested. A key piece of
that structure is a \emph{Frobenius monoid} on each object 
(see Appendix~\ref{appendix:frobenius}): a monoid $(\mu, \eta)$ 
and a comonoid $(\delta, \varepsilon)$ on the same object, satisfying the Frobenius law
\[
	(\mathrm{id} \otimes \mu) \circ (\delta \otimes \mathrm{id})
	\;=\;
	\delta \circ \mu
	\;=\;
	(\mu \otimes \mathrm{id}) \circ (\mathrm{id} \otimes \delta).
\]
For a hypergraph category, this structure must also be \emph{commutative}
and \emph{special}: the multiplication and comultiplication respect the symmetry,
and $\mu\circ\delta=\mathrm{id}$. A symmetric monoidal category with a chosen
special commutative Frobenius monoid on every object, compatibly with the tensor
product and its unit, is called a \emph{hypergraph category}
\cite[Definition~2.12]{fong2019hypergraph}.

Hypergraph categories carry a self-dual compact closed structure
\cite[Proposition~3.1]{fong2019hypergraph}. For each object $A$, its Frobenius
operations define a cup and a cap by
\[
	\mathrm{cup}_A := \delta_A\circ\eta_A : I\to A\otimes A,
	\qquad
	\mathrm{cap}_A := \varepsilon_A\circ\mu_A : A\otimes A\to I.
\]
These satisfy the snake equations, so $A$ can serve as its own dual.
It is this richer structure that we shall adopt as our ambient framework.

\begin{remark}
	If $Q$ is commutative, then $\mathbf{Rel}(Q)$ has a canonical special
commutative Frobenius structure on every object: its maps copy, discard,
and take the relational converses of those operations, with weight $e$ on
allowed pairs and $\bot$ otherwise~\cite{coecke2018generalizedrelations}.
\end{remark}

This structure will justify identities in the calculi developed below.
To illustrate its operations, we introduce the following symbols.

\begin{figure}[H]
    \centering
    \begin{center}
    \begin{tikzpicture}
        \pic (g1) at (0.500,0.500) {copy};
        \node[scale=1.5] at (1.25,0.5) {$|$};
        \pic (g2) at (2.000,0.500) {discard};
        \node[scale=1.5] at (3,0.5) {$|$};
        \pic (g3) at (3.500,0.500) {codiscard};
        \node[scale=1.5] at (4.25,0.5) {$|$};
        \pic (g4) at (5.000,0.500) {cocopy};
        \node at (0.5,-0.5) {\footnotesize copy};
        \node at (2.0,-0.5) {\footnotesize discard};
        \node at (3.5,-0.5) {\footnotesize unit};
        \node at (5.0,-0.5) {\footnotesize comp};
    \end{tikzpicture}

    \vspace{0.5em}

    \begin{tikzpicture}
        \pic (g1) at (0.500,0.500) {comultiplication};
        \node[scale=1.5] at (1.25,0.5) {$|$};
        \pic (g2) at (2.000,0.500) {cozero};
        \node[scale=1.5] at (3,0.5) {$|$};
        \pic (g3) at (3.500,0.500) {zero};
        \node[scale=1.5] at (4.25,0.5) {$|$};
        \pic (g4) at (5.000,0.500) {multiplication};
        \node at (0.5,-0.5) {\footnotesize coadd};
        \node at (2.0,-0.5) {\footnotesize cozero};
        \node at (3.5,-0.5) {\footnotesize zero};
        \node at (5.0,-0.5) {\footnotesize add};
    \end{tikzpicture}
\end{center}

    \caption{Frobenius additive(black) and co-additive(white) structure}
\end{figure}

The black structure consists of copying and discarding operations and their
relational converses. These converses enforce equality between two values or
create an arbitrary value, respectively. The copying and additive structures
on linear relations are described in~\cite[Examples~2.10--2.11]{fong2019hypergraph}.
The Frobenius law gives the identities below. Copying the first value and
comparing it with the second is equivalent to copying the second and comparing
it with the first, or to comparing both values and then copying the common value.

\begin{figure}[H]
    \centering
    \begin{center}
    \begin{tikzpicture}[axpic]
      \begin{scope}[shift={(0.000,1.000)}]
      \pic (g1) at (0.500,0.500) {copy};
      \end{scope}
      \coordinate (id_in_2) at (0,0.000);
      \coordinate (id_out_3) at (1.000,0.000);
      \draw (id_in_2) -- (id_out_3);
      \begin{scope}[shift={(1.250,0.000)}]
      \begin{scope}[shift={(0.000,1.000)}]
      \coordinate (id_in_4) at (0,1.000);
      \coordinate (id_out_5) at (1.000,1.000);
      \draw (id_in_4) -- (id_out_5);
      \end{scope}
      \pic (g6) at (0.500,0.500) {cocopy};
      \end{scope}
      \draw (g1-out-0) to[out=0,in=180] (id_in_4);
      \draw (g1-out-1) to[out=0,in=180] (g6-in-0);
      \draw (id_out_3) to[out=0,in=180] (g6-in-1);
    \end{tikzpicture}
    \;\begin{tikzpicture}[axpic]\node{$=$};\end{tikzpicture}\;
    \begin{tikzpicture}[axpic]
      \pic (g1) at (0.500,0.500) {cocopy};
      \begin{scope}[shift={(1.250,0.000)}]
      \pic (g2) at (0.500,0.500) {copy};
      \end{scope}
      \draw (g1-out-0) to[out=0,in=180] (g2-in-0);
    \end{tikzpicture}
    \;\begin{tikzpicture}[axpic]\node{$=$};\end{tikzpicture}\;
    \begin{tikzpicture}[axpic]
      \begin{scope}[shift={(0.000,1.000)}]
      \coordinate (id_in_1) at (0,1.000);
      \coordinate (id_out_2) at (1.000,1.000);
      \draw (id_in_1) -- (id_out_2);
      \end{scope}
      \pic (g3) at (0.500,0.500) {copy};
      \begin{scope}[shift={(1.250,0.000)}]
      \begin{scope}[shift={(0.000,1.000)}]
      \pic (g4) at (0.500,0.500) {cocopy};
      \end{scope}
      \coordinate (id_in_5) at (0,0.000);
      \coordinate (id_out_6) at (1.000,0.000);
      \draw (id_in_5) -- (id_out_6);
      \end{scope}
      \draw (id_out_2) to[out=0,in=180] (g4-in-0);
      \draw (g3-out-0) to[out=0,in=180] (g4-in-1);
      \draw (g3-out-1) to[out=0,in=180] (id_in_5);
    \end{tikzpicture}
\end{center}

\end{figure}

\section{Graphical Linear Algebra}\label{sec:linear-algebra}

We now introduce an equational, diagrammatic syntax for linear algebra.
We recall the relevant results without reproducing their proofs. The linear
presentation and its proofs are given in~\cite{bonchi2017interacting}, and the
affine extension is developed in~\cite{bonchi2019affine}.
The categorical and graphical background is developed in
Appendix~\ref{appendix:category-theory} and Section~\ref{sec:string-diagrams}.

Fix a field $\mathbb K$; all applications in this thesis use
$\mathbb K=\mathbb R$.

\begin{definition}[Linear relation]\label{def:linear-relations}
A \emph{linear relation} $R:m\to n$ is a linear subspace
\[
 R\subseteq\mathbb K^m\oplus\mathbb K^n.
\]
We write $x\mathrel{R}y$ when $(x,y)\in R$.  The PROP
$\mathbf{LinRel}_{\mathbb K}$ has natural numbers as objects and linear
relations as morphisms.  Its identity, composition, and tensor are
\begin{align*}
 \mathrm{id}_n
   &=\{(x,x)\mid x\in\mathbb K^n\},\\
 S\circ R
   &=\{(x,z)\mid \exists y\in\mathbb K^n:
                    x\mathrel{R}y\text{ and }y\mathrel{S}z\},\\
 R\otimes S
   &=\{((x,u),(y,v))\mid x\mathrel{R}y,
                              \ u\mathrel{S}v\}.
\end{align*}
After the canonical identification
$\mathbb K^m\oplus\mathbb K^p\cong\mathbb K^{m+p}$, the tensor on objects is
addition.
\end{definition}

The composite is again a linear subspace: the simultaneous constraints
form a homogeneous linear system, and projecting away the intermediate
variables preserves linearity.
\footnote{Composition is relational composition, computed by projecting the
    simultaneous linear constraints onto the external variables.}

Every matrix $A\in\mathbb K^{n\times m}$ determines the functional linear
relation
\[
 \operatorname{graph}(A)
 =\{(x,y)\in\mathbb K^m\oplus\mathbb K^n\mid y=Ax\}.
\]
Connecting these functional relations recovers ordinary matrix multiplication
over $\mathbb K$:
\[
 \operatorname{graph}(B)\circ\operatorname{graph}(A)
 =\operatorname{graph}(BA).
\]
Graphical Linear Algebra is a graphical presentation of
$\mathbf{LinRel}_{\mathbb K}$~\cite{bonchi2017interacting}. Its diagrams are
generated by two interacting families of nodes and one scalar box for each
$k\in\mathbb K$.
Reflecting a diagram exchanges its input and output interfaces. If the original
diagram denotes $R:m\to n$, its reflection denotes the converse relation
$R^{\mathrm{op}}:n\to m$, defined by
\[
 x\mathrel{R}y \iff y\mathrel{R^{\mathrm{op}}}x.
\]
Reflection is a contravariant involution: it reverses the order of composition,
and reflecting twice recovers the original relation. In general, $R$ and
$R^{\mathrm{op}}$ are different typed morphisms; reflection does not assert
that they are equal or that $R$ is invertible.

The following table gives the type and exact relation denoted
by every basic orientation; the reversed orientation denotes the converse
relation.

\begin{table}[H]
\centering
\caption{Generators of Graphical Linear Algebra and their linear-relation semantics.}
\label{tab:appendix-gla-generators}
\begin{tabular}{c c p{8.1cm}}
\\
\toprule
\textbf{diagram} & \textbf{type} & \textbf{interpretation} \\
\midrule
\begin{tikzpicture}[axpic]\pic (g) at (0,0) {multiplication};\end{tikzpicture}
& $2\to1$ & $\{((x_1,x_2),y)\mid y=x_1+x_2\}$ \\
\begin{tikzpicture}[axpic]\pic (g) at (0,0) {zero};\end{tikzpicture}
& $0\to1$ & $\{(*,0)\}$ \\
\begin{tikzpicture}[axpic]\pic (g) at (0,0) {copy};\end{tikzpicture}
& $1\to2$ & $\{(x,(y_1,y_2))\mid y_1=x=y_2\}$ \\
\begin{tikzpicture}[axpic]\pic (g) at (0,0) {discard};\end{tikzpicture}
& $1\to0$ & $\{(x,*)\mid x\in\mathbb K\}$ \\
\begin{tikzpicture}[axpic]\pic[val=k] (g) at (0,0) {scalar};\end{tikzpicture}
& $1\to1$ & $\{(x,y)\mid y=kx\}$, for $k\in\mathbb K$ \\
\bottomrule
\end{tabular}
\end{table}

The white node is addition and the black node is equality.  Their reflected
forms have the following readings:
\begin{itemize}
    \item white comultiplication $1\to2$ imposes $x=y_1+y_2$;
    \item white counit $1\to0$ imposes $x=0$;
    \item black multiplication $2\to1$ imposes $x_1=x_2=y$;
    \item black unit $0\to1$ allows an arbitrary value of $y$;
    \item a reflected scalar imposes $x=ky$.
\end{itemize}
Thus direction distinguishes solving an equation forward from treating the
same equation as a constraint.  When $k\ne0$, a reflected $k$-box is equivalent
to a forward $k^{-1}$-box; when $k=0$, its semantics is the non-functional
constraint $x=0$ with arbitrary output.

These generators allow us to represent an arbitrary matrix:

\begin{definition}[Matrix diagram]
Let $A=(a_{ij})\in\mathbb K^{n\times m}$. The matrix diagram
$A:m\to n$ is built from the generators above as follows. Copy each input
$x_j$ into $n$ wires, one for each output. On the wire from input $j$ to
output $i$, place the scalar box $a_{ij}$. Reorder the wires to group them
by output, and add the $m$ contributions in each group. Thus output $i$ is
\[
 y_i=\sum_{j=1}^{m}a_{ij}x_j,
\]
and the diagram denotes exactly $\operatorname{graph}(A)$.

The following example illustrates the construction for a $2\times2$ matrix;
the same procedure applies to an $n\times m$ matrix.

\centering
\diagdef{
    \pic[val=A] {matrix};
}{
              \begin{scope}[shift={(0.000,1.000)}]
              \begin{scope}[shift={(0.000,1.000)}]
              \pic (g1) at (0.500,0.500) {copy};
              \end{scope}
              \pic (g2) at (0.500,0.500) {copy};
              \end{scope}
              \begin{scope}[shift={(1.250,0.000)}]
              \begin{scope}[shift={(0.000,1.000)}]
              \begin{scope}[shift={(0.000,1.000)}]
              \begin{scope}[shift={(0.000,1.000)}]
                  \pic (g3)[val=a_{1,1}] at (0.500,0.500) {scalar};
              \end{scope}
              \pic (g4)[val=a_{2,1}] at (0.500,0.500) {scalar};
              \end{scope}
              \pic (g5)[val=a_{1,2}] at (0.500,0.500) {scalar};
              \end{scope}
              \pic (g6)[val=a_{2,2}] at (0.500,0.500) {scalar};
              \end{scope}
              \draw (g1-out-0) to[out=0,in=180] (g3-in-0);
              \draw (g1-out-1) to[out=0,in=180] (g4-in-0);
              \draw (g2-out-0) to[out=0,in=180] (g5-in-0);
              \draw (g2-out-1) to[out=0,in=180] (g6-in-0);
              \begin{scope}[shift={(2.500,1.000)}]
              \begin{scope}[shift={(0.000,1.000)}]
              \pic (g7) at (0.500,0.500) {multiplication};
              \end{scope}
              \pic (g8) at (0.500,0.500) {multiplication};
              \end{scope}
              \draw (g3-out-0) to[out=0,in=180] (g7-in-0);
              \draw (g4-out-0) to[out=0,in=180] (g7-in-1);
              \draw (g5-out-0) to[out=0,in=180] (g8-in-0);
              \draw (g6-out-0) to[out=0,in=180] (g8-in-1);
}


\end{definition}

The equational theory of Graphical Linear Algebra can be organized into four layers.

\begin{enumerate}[label=\textnormal{(\roman*)}]
    \item The addition node and zero form a commutative monoid, while copy and
    discard form a cocommutative comonoid.  Reflecting these equations gives
    the corresponding structures in the opposite direction.

    \item Each color, together with its reflection, is a special commutative
    Frobenius algebra.  The spider theorem therefore applies separately to
    connected monochromatic regions.

    \item The two colors interact by the bialgebra and Hopf laws.  These laws
    express that copying preserves addition and zero, and that additive
    inverses cancel.  The scalar $-1$ is the antipode.

    \item Scalar boxes obey the arithmetic of $\mathbb K$: consecutive boxes
    multiply, parallel copied-and-added boxes add, $1$ is the identity, and
    scalars commute with the appropriate copying and addition nodes.

\end{enumerate}

Every equation is sound because both sides denote the same linear subspace.
For example, copying after addition and adding copied components both denote
\[
 \{((x_1,x_2),(y_1,y_2))\mid
       y_1=x_1+x_2=y_2\}.
\]
The equations do more than preserve semantics: they characterize it.

\begin{definition}[The linear syntax]\label{def:linear-syntax}
We call $\mathbf{Graphical\ Linear\ Algebra}_{\mathbb K}$ the PROP built from
 the linear generators and equations presented above~\cite{bonchi2017interacting}.
 A diagram $m\to n$ has $m$ input wires and $n$ output wires. We build
 diagrams by connecting these generators or placing them in parallel,
 and use the equations to identify them, as in our definition of a
 graphical theory. The affine point and its axioms are not included.
\end{definition}

Notice that we already have a language here.
We can express $x+y=0$, but not $x+y=1$. Why? Because every linear relation
admits the zero assignment, and the second equation does not. The theorem
below tells us that this smaller language captures exactly
$\mathbf{LinRel}_{\mathbb K}$. We will start from this same language when
we add inequalities in Chapter~\ref{sec:gpa}.

\begin{theorem}[Presentation theorem for Graphical Linear Algebra]\label{thm:linear-algebra-presentation}
For a field $\mathbb K$, the standard interpretation extends to an isomorphism
of PROPs
\[
 \llbracket{-}\rrbracket:
 \mathbf{Graphical\ Linear\ Algebra}_{\mathbb K}\xrightarrow{\ \cong\ }
 \mathbf{LinRel}_{\mathbb K}.
\]
Consequently, Graphical Linear Algebra is sound, complete, and expressive for linear
relations~\cite{bonchi2017interacting}.
\end{theorem}

Linear relations are more general than linear maps. A reflected matrix
denotes the converse relation
\[
 \operatorname{graph}(A)^{\mathrm{op}}=\{(y,x)\mid y=Ax\},
\]
which may be partial or multivalued. Every linear relation has a homogeneous
cospan form $Ex=Fy$, and hence a diagrammatic representative.

A matrix box $A:m\to n$ abbreviates any Graphical Linear Algebra diagram whose interpretation is
$y=Ax$.  Its reflected box denotes the converse relation.  For compatible
matrices, connecting boxes gives matrix multiplication, while copying inputs,
applying two matrices, and adding their outputs gives the corresponding block
matrix.  A general linear relation box may be replaced by its homogeneous cospan form
\[
 x\mathrel{R}y\quad\Longleftrightarrow\quad Ex=Fy.
\]

\section{Graphical Affine Algebra}

Without the affine point, all generators denote homogeneous linear relations.
We now extend Graphical Linear Algebra with the affine point $\rtop:0\to1$, which
denotes the constant $1$. Together with scalar multiplication and copying,
it constructs constants and translations. The resulting language is
Graphical Affine Algebra~\cite{bonchi2019affine}.

\[
\begin{tikzpicture}[baseline=(current bounding box.center)]
    \node (lbr) at (0,0) {$\llbracket$};

    \pic (g) at (0.3,0) {affine};

    \node (rbr) at (1.1,0) {$\rrbracket$};
    \node (eq)  at (1.5,0) {$=$};
    \node (rhs) at (2.4,0) {$\{(*,1)\}$};
\end{tikzpicture}
\]
\begin{definition}[The affine syntax]\label{def:affine-syntax}
We call $\mathbf{Graphical\ Affine\ Algebra}_{\mathbb K}$ the PROP built by adding
 the affine point and its three equations to the linear syntax of
 Definition~\ref{def:linear-syntax}. We build and identify diagrams in
 the same way as before, now using both the linear and affine equations.
\end{definition}

What changes is that we can now supply the constant $1$. Multiplying it
by a scalar gives any constant we need, and adding it to a variable lets
us write equations such as $x+y=1$. More generally, $Ex=Fy$ becomes
$Ex+b=Fy$. We have not changed how diagrams connect; we have added one
component and the equations that tell us how to use it. If we leave that
component out, we are back to diagrams describing linear relations.

The same distinction will matter in the next chapter. Adding inequalities
to the smaller linear language gives polyhedral cones. Adding them to the
affine language gives general polyhedra. Once again, allowing constants
is what makes the difference.

This extension illustrates the incremental construction of the language:
a new generator and its equations add a specific semantic capability.
The same approach will be used to introduce inequalities and objectives,
extending the calculus to polyhedra and linear programs. Comparing these
presentations identifies the additional operations used to pass from one
class of relations to another.


\begin{definition}[Affine relation]
An \emph{affine relation} $R:m\to n$ is an affine subspace (or the empty set)
\[
 R\subseteq\mathbb K^m\oplus\mathbb K^n.
\]
We write $x\mathrel{R}y$ when $(x,y)\in R$.  The PROP
$\mathbf{AffRel}_{\mathbb K}$ has natural numbers as objects and affine
relations as morphisms.  Its identity, composition, and tensor are
\begin{align*}
 \mathrm{id}_n
   &=\{(x,x)\mid x\in\mathbb K^n\},\\
 S\circ R
   &=\{(x,z)\mid \exists y\in\mathbb K^n:
                    x\mathrel{R}y\text{ and }y\mathrel{S}z\},\\
 R\otimes S
   &=\{((x,u),(y,v))\mid x\mathrel{R}y,
                              \ u\mathrel{S}v\}.
\end{align*}
After the canonical identification
$\mathbb K^m\oplus\mathbb K^p\cong\mathbb K^{m+p}$, the tensor on objects is
addition.
\end{definition}

The composite is again an affine subspace or empty: the simultaneous
constraints $(x,y)\in R$ and $(y,z)\in S$ form an affine system, and projecting
away $y$ preserves affinity.  Hence composition of affine relations is exactly
affine variable elimination.

Affine relations are nevertheless strictly more general than affine maps.
Linear relations embed as the special case with zero translation. More generally,
every affine relation admits a presentation by affine linear equations:
\[
 R=\{(x,y)\mid Ex+b=Fy\}
\]
for suitable matrices $E,F$ and vector $b$. This affine cospan form is the
semantic origin of Graphical Affine Algebra normal forms.

\begin{remark}
    Every morphism in $\textbf{AffRel}$ admits the following diagram representation
    \begin{figure}[h]
    \centering
    \begin{tikzpicture}
          \pic[val=E] (g1) at (0.500,1.000) {matrix};
          \pic (g3) at (2.500,0.500) {multiplication};
          \pic (g4) at (0.000,-0.000) {affine};
          \pic[val=b] (g5) at (0.750,0.000) {scalar};
          \pic[val=F] (g2) at (3.750,0.500) {comatrix};
          \draw (g1-out-0) to[out=0,in=180] (g3-in-0);
          \draw (g4-out-0) to[out=0,in=180] (g5-in-0);
          \draw (g5-out-0) to[out=0,in=180] (g3-in-1);
          \draw (g3-out-0) to[out=0,in=180] (g2-in-0);
    \end{tikzpicture}
    \end{figure}
\end{remark}

\begin{theorem}[Presentation theorem for Graphical Affine Algebra]\label{thm:appendix-gla-presentation}
For a field $\mathbb K$, the standard interpretation extends to an isomorphism
of PROPs
\[
 \llbracket{-}\rrbracket:
 \mathbf{Graphical\ Affine\ Algebra}_{\mathbb K}\xrightarrow{\ \cong\ }
 \mathbf{AffRel}_{\mathbb K}.
\]
Consequently, Graphical Affine Algebra is sound, complete, and expressive for affine
relations~\cite{bonchi2019affine}.
\end{theorem}

The theorem explains why Graphical Affine Algebra is useful in the main text.  A diagram is not
merely a picture of a chosen matrix calculation: it denotes the entire
solution relation of an affine system.  Sequential composition existentially
eliminates internal variables, tensor combines independent systems, and
reflection exchanges a relation with its converse.  Diagram rewriting is
therefore a complete algebra of affine systems.

The particular derived equalities invoked by the proofs in
Chapters~\ref{sec:gpa} and~\ref{sec:glp} are collected below.  They are grouped
as (G1) matrix calculus, (G2) transposition and cospan form, (G3) bending, and
(G4) the black spider laws.  All follow from the Graphical Linear Algebra and Graphical Affine Algebra equations summarized above and
are included here so that later diagrammatic derivations can be checked without
expanding every matrix box into scalar generators.

% Rules of the underlying linear and affine algebras (GLA, GAA) that are used
% as background throughout, but are not contributions of this thesis.
\medskip
\noindent
\fbox{\begin{minipage}{0.97\linewidth}
    \begin{center}\textbf{(G1) Matrices}\quad{\footnotesize\cite{bonchi2017interacting,paixao2022highlevel}}\end{center}
    \begin{center}\begin{tabular}{r c c c}
      \diagrow{(comp.)}{
          \pic[val=A] (g1) at (0.250,0.500) {matrix};
          \pic[val=B] (g2) at (2.250,0.500) {matrix};
          \draw (g1-out-0) to[out=0,in=180] (g2-in-0);
      }{
          \pic[val=BA] (g1) at (0.500,0.500) {matrix};
      }
      \diagrow{(row)}{
          \pic[val=A] (g1) at (0.500,1.000) {matrix};
          \pic[val=B] (g2) at (0.500,0.000) {matrix};
          \pic (g3) at (2.050,0.500) {multiplication};
          \draw (g1-out-0) to[out=0,in=180] (g3-in-0);
          \draw (g2-out-0) to[out=0,in=180] (g3-in-1);
      }{
          \pic[val={[A \ B]}] (g1) at (1.000,0.000) {2dmatrix};
          \draw (0.000,0.300) -- (g1-in-0);
          \draw (0.000,-0.300) -- (g1-in-1);
      }
      \diagrow{(antipode)}{
          \pic (g1) at (0.500,0.500) {antipode};
      }{
          \pic[val=-1] (g1) at (0.500,0.500) {scalar};
      }
      \diagrow{(zero sc.)}{
          \pic[val=0] (g1) at (0.500,0.500) {scalar};
      }{
          \pic (g1) at (0.500,0.500) {discard};
          \pic (g2) at (1.750,0.500) {zero};
          \draw (0.000,0.500) -- (g1-in-0);
          \draw (g2-out-0) -- (2.500,0.500);
      }
      \diagrow{(inversion)}{
          \pic[val=k] (g1) at (0.500,0.500) {scalar};
      }{
          \pic[val=k^{-1}] (g1) at (0.500,0.500) {coscalar};
      }
    \end{tabular}\end{center}
\end{minipage}}

\medskip
\noindent
\fbox{\begin{minipage}{0.97\linewidth}
    \begin{center}\textbf{(G2) Transposition and cospan form}\quad{\footnotesize\cite{bonchi2017interacting}}\end{center}
    A diagram may be slid along a cap, exchanging it for its transpose; on
    $\mathbf{GLA}$ diagrams the transpose is the converse relation, so a matrix
    becomes a comatrix.
    \begin{center}\begin{tabular}{r c c c}
      \diagrow{(slide)}{
          \pic[val=A] (g1) at (0.500,1.000) {matrix};
          \draw (g1-out-0) -- (1.600,1.000) arc (90:-90:0.500) -- (0.000,0.000);
      }{
          \pic[val=A] (g1) at (0.500,0.000) {comatrix};
          \draw (g1-out-0) -- (1.600,0.000) arc (-90:90:0.500) -- (0.000,1.000);
      }
      \diagrow{(cospan)}{
          \pic[val=g] (g1) at (0.500,0.500) {relation};
      }{
          \pic[val=E] (g1) at (0.500,1.000) {matrix};
          \pic (g3) at (2.500,0.500) {multiplication};
          \pic (g4) at (0.000,-0.000) {affine};
          \pic[val=b] (g5) at (0.750,0.000) {scalar};
          \draw (g1-out-0) to[out=0,in=180] (g3-in-0);
          \draw (g4-out-0) to[out=0,in=180] (g5-in-0);
          \draw (g5-out-0) to[out=0,in=180] (g3-in-1);
      }
    \end{tabular}\end{center}
\end{minipage}}

\medskip
\noindent
\fbox{\begin{minipage}{0.97\linewidth}
    \begin{center}\textbf{(G3) Bending}\quad{\footnotesize\cite{bonchi2017interacting,bonchi2019affine}}\end{center}
    The white cap is the black cap with an antipode on one leg; this is what
    turns an equation into a pair of bendable wires.
    \begin{center}\begin{tabular}{r c c c}
      \diagrow{(white cap)}{
          \pic (g1) at (0.500,0.500) {multiplication};
          \pic (g2) at (1.750,0.500) {cozero};
          \draw (g1-out-0) to[out=0,in=180] (g2-in-0);
      }{
          \pic (g1) at (0.500,0.000) {antipode};
          \draw (g1-out-0) -- (1.600,0.000) arc (-90:90:0.500) -- (0.000,1.000);
      }
      \diagrow{(unbend $\oplus$)}{
          \pic (g1) at (0.500,0.500) {comultiplication};
          \pic (g2) at (1.750,0.000) {antipode};
          \draw (g1-out-0) -- (2.500,1.000);
          \draw (g1-out-1) to[out=0,in=180] (g2-in-0);
          \draw (g2-out-0) -- (2.500,0.000);
      }{
          \pic (g1) at (1.000,0.500) {multiplication};
          \draw (0.000,1.000) to[out=0,in=180] (g1-in-0);
          \draw (g1-out-0) -- (1.900,0.500) arc (90:-90:0.500) -- (0.000,-0.500);
      }
    \end{tabular}\end{center}
\end{minipage}}

\medskip
\noindent
\fbox{\begin{minipage}{0.97\linewidth}
    \begin{center}\textbf{(G4) The black structure is a spider}\quad{\footnotesize\cite{bonchi2017interacting}}\end{center}
    Any connected network of black nodes with $a$ inputs and $b$ outputs equals
    the single node $\mathsf{cc}_a;\mathsf{cp}_b$; the two laws below are the
    cases used here.
    \begin{center}\begin{tabular}{r c c c}
      \diagrow{(special)}{
          \pic (g1) at (0.500,0.500) {copy};
          \pic (g2) at (1.750,0.500) {cocopy};
          \draw (g1-out-0) to[out=0,in=180] (g2-in-0);
          \draw (g1-out-1) to[out=0,in=180] (g2-in-1);
      }{
          \draw (0.000,0.500) -- (1.000,0.500);
      }
      \diagrow{(merge/del)}{
          \pic (g1) at (0.500,0.500) {cocopy};
          \pic (g2) at (1.750,0.500) {discard};
          \draw (g1-out-0) -- (g2-in-0);
      }{
          \pic (g1) at (0.500,1.000) {discard};
          \pic (g2) at (0.500,0.000) {discard};
          \draw (0.000,1.000) -- (g1-in-0);
          \draw (0.000,0.000) -- (g2-in-0);
      }
      \diagrow{(split/un)}{
          \pic (g1) at (0.500,0.500) {codiscard};
          \pic (g2) at (1.750,0.500) {copy};
          \draw (g1-out-0) -- (g2-in-0);
      }{
          \pic (g1) at (0.500,1.000) {codiscard};
          \pic (g2) at (0.500,0.000) {codiscard};
          \draw (g1-out-0) -- (1.250,1.000);
          \draw (g2-out-0) -- (1.250,0.000);
      }
    \end{tabular}\end{center}
\end{minipage}}


The affine point additionally turns homogeneous matrix equations into affine
ones: together with copying and scalar multiplication it constructs any
constant vector $b$, and hence any constraint $Ex+b=Fy$.  The later calculi in
this thesis extend this affine base with order and optimization generators.
Their new equations are studied in the later chapters, whereas the affine
rules on which they rely are established here. The categorical background
is developed in Appendix~\ref{appendix:category-theory}.
